% ============================================================
%  Problem Formulation and Motivation
% ============================================================
\section{Problem and Method}
\label{sec:problem}

\subsection{Task Definition}
Let $f_I$ and $f_T$ be frozen image and text encoders and $u=f_I(I)$, $v=f_T(T)$.
The raw score is the cosine similarity $s(I,T)=\operatorname{sim}(u,v)$, and each
pair carries a binary label $y\in\{0,1\}$ (clean or noisy). The task is to construct
$q(I,T)$ whose ranking separates the classes better than $s$. Two properties
distinguish it from same-image candidate ranking: the comparison is \emph{global}
(the two pairs share neither image nor caption) and scoring is \emph{pair-local}
($q$ may not depend on which other pairs are present). We therefore report cross-pair
AUROC and cycle-level 2AFC (two-alternative forced choice) separately.

\subsection{Empirical Limitation of Raw Similarity}
Figure~\ref{fig:motivation} shows a COCO Quantity example scored by one frozen CLIP
ViT-B/16: the correct pair scores $0.252$, the pair whose caption misstates the
number of animals $0.331$. A single case proves little, so we quantify it on the
largest facet where the effect is measurable, Visual Genome Attribute. Comparing every clean
pair with every noisy pair, $31.9\%$ of the comparisons place the noisy pair above
the clean one, which equals $1-\mathrm{AUROC}$ ($1-0.6806$); the corresponding value
for Quantity is $16.9\%$ ($1-0.8305$). Within one reciprocal cycle, replacing the
caption of a fixed image changes the score ($0.324\to0.188$) and replacing the image
changes it too ($0.324\to0.253$), so the score depends on which image and caption are
involved, not solely on whether they correspond. These are conditioned variations,
not measurements of an isolated modality bias.

\subsection{Hypothesis: Additive Score Components}
Suppose the standardised score decomposes as
\begin{equation}
  z \;=\; \underbrace{m(I,T)}_{\text{cross-modal match}} \;+\;
         \underbrace{c(I,T)}_{\text{additive component}},
  \label{eq:decomp}
\end{equation}
where $c$ is predictable from $u$ and $v$ and largely independent of whether they
correspond. Under \eqref{eq:decomp}, a stronger encoder raises the first term while
the ordering of unrelated pairs remains governed by the second, consistent with the
unchanged Attribute AUROC of ViT-L/14. The decomposition also explains why a monotone
recalibration cannot help: it acts on $z$ as a whole. Equation~\eqref{eq:decomp} is a
modelling assumption, not a generative account, and we do not claim that $c$ is
identifiable as image or text noise. What is tested is its operative consequence: if
a component of $z$ is predictable from $u$ and $v$ alone, removing that prediction
should improve global discrimination, and a control fitted to permuted targets
should not.
